% Add competences with simplified syntax
%Chapitre 1 : Calcul numérique
\AddCompetence{CalcNum1}{1M}{1}{Calculer avec les nombres rationnels}
\AddCompetence{CalcNum2}{1M}{1}{Utiliser la division euclidienne}
%Chapitre 2 : Ensembles
\AddCompetence{Ens1}{1M}{2}{Connaître les ensembles de nombres}
%Chapitre 3 : Calcul littéral
\AddCompetence{CalcLit0}{1M}{3}{Connaître le vocabulaire du calcul littéral}
\AddCompetence{CalcLit1}{1M}{3}{Additionner\, soustraire\, multiplier des monômes}
\AddCompetence{CalcLit2}{1M}{3}{Additionner\, soustraire des polynômes}
\AddCompetence{CalcLit3}{1M}{3}{Réduire une expression littérale}
\AddCompetence{CalcLit4}{1M}{3}{Appliquer la distributivité simple}
\AddCompetence{CalcLit5}{1M}{3}{Appliquer la distributivité double}
\AddCompetence{CalcLit6}{1M}{3}{Utiliser le calcul littéral pour modéliser une situation}
\AddCompetence{CalcLit7}{1M}{3}{Évaluer une expression littérale}
\AddCompetence{IdRem}{1M}{3}{Développer une expression en utilisant l'identité remarquable}
\AddCompetence{Fact1}{1M}{3}{Factoriser la mise en évidence}
\AddCompetence{Fact2}{1M}{3}{Factoriser à l'aide des identités remarquables}
\AddCompetence{Fact3}{1M}{3}{Factoriser des groupements}
\AddCompetence{Fact4}{1M}{3}{Factoriser par enchaînement des techniques}
\AddCompetence{CalcLit7}{1M}{3}{Factoriser à l'aide de la mise en évidence}
%Chapitre 4 : Équations du premier degré
\AddCompetence{Eq1}{1M}{4}{Résoudre une équation du premier degré dont les coefficient sont rationnels}
\AddCompetence{Eq1Irr}{1M}{4}{Résoudre une équation du premier degré dont les coefficient sont irrationnels}
\AddCompetence{PE1}{1M}{4}{Réciter et appliquer le premier principe d'équivalence}
\AddCompetence{PE2}{1M}{4}{Réciter et appliquer le deuxième principe d'équivalence}
\AddCompetence{ProbEq1}{1M}{4}{Résoudre une situation problème par une équation du premier degré}
\AddCompetence{EqParam}{1M}{4}{Résoudre une équation du premier degré à un paramètre}
%Chapitre 5 : Équations du second degré
\AddCompetence{ProdNul}{1M}{5}{Résoudre une équation du second degré en appliquant le théorème du produit nul}
\AddCompetence{ResFact}{1M}{5}{Résoudre une équation en utilisant les méthodes de factorisation}
\AddCompetence{ComplCarre}{1M}{5}{Résoudre une équation à l'aide de la méthode de complétion du carré}
\AddCompetence{FactTri}{1M}{5}{Factoriser un trinôme en appliquant la formule du deuxième degré}
\AddCompetence{ResDelta}{1M}{5}{Résoudre une équation du second degré en appliquant la formule du deuxième degré ($\Delta$)}
\AddCompetence{EqIrr}{1M}{5}{Résoudre une équation irrationnelle}
\AddCompetence{Bicarr}{1M}{5}{Résoudre une équation bicarrée}
\AddCompetence{ProbEq2}{1M}{5}{Résoudre une situation problème par une équation du second degré}
%Chapitre 6 : Systemes d'équations
\AddCompetence{PE3}{1M}{6}{Réciter et appliquer le troisième principe d'équivalence}
\AddCompetence{CombLin}{1M}{6}{Résoudre un système d'équations par combinaison linéaire}
\AddCompetence{CombSub}{1M}{6}{Résoudre un système d'équations par substitution}
\AddCompetence{3f3}{1M}{6}{Résoudre un système 3x3 par substitution et combinaison linéaire}
\AddCompetence{ProbSys}{1M}{6}{Résoudre une situation problème par un système d'équations}
%Chapitre 7 : Fonctions
\AddCompetence{Fonc1}{1M}{7}{Connaître le vocabulaire des fonctions}
\AddCompetence{Fonc2}{1M}{7}{Déterminer l'image d'un nombre par une fonction}
\AddCompetence{Fonc3}{1M}{7}{Déterminer la préimage d'un nombre par une fonction}
\AddCompetence{Fonc4}{1M}{7}{Déterminer les zéros d'une fonction}
\AddCompetence{Fonc5}{1M}{7}{Établir et compléter un tableau de valeurs}
\AddCompetence{Fonc6}{1M}{7}{Déterminer si une fonction passe par un point donné en évaluant la fonction}
\AddCompetence{Fonc7}{1M}{7}{Déterminer la pente et l'ordonnée à l'origine d'une droite à partir de son expression algébrique}
\AddCompetence{Fonc8}{1M}{7}{Représenter graphiquement une fonction affine}
\AddCompetence{Fonc9}{1M}{7}{Distinguer les différentes écritures algébriques d'une droite~: équation cartésienne, équation réduite}
\AddCompetence{Fonc10}{1M}{7}{Déterminer l'expression algébrique d'une droite à partir de sa représentation graphique}
\AddCompetence{Fonc11}{1M}{7}{Déterminer l'expression algébrique d'une droite à partir de deux points}
\AddCompetence{Fonc12}{1M}{7}{Connaître les conditions de parallélisme et de perpendicularité de deux droites}
\AddCompetence{Fonc13}{1M}{7}{Déterminer l'expression algébrique d'une droite à partir de conditions données sur la pente et un point}
\AddCompetence{Fonc14}{1M}{7}{Déterminer si trois points sont alignés}
\AddCompetence{Fonc15}{1M}{7}{Déterminer le point d'intersection de deux droites}
\AddCompetence{Fonc16}{1M}{7}{Utiliser les différentes écritures d'une parabole~: forme générale, forme canonique, forme factorisée}
\AddCompetence{Fonc17}{1M}{7}{Établir un tableau des signes d'une parabole}
\AddCompetence{Fonc18}{1M}{7}{Réaliser une étude complète d'une parabole}
\AddCompetence{Fonc19}{1M}{7}{Déterminer le(s) point(s) d'intersection d'une droite et d'une parabole}
\AddCompetence{Fonc20}{1M}{7}{Déterminer le(s) point(s) d'intersection de deux paraboles}
\AddCompetence{Fonc20}{1M}{7}{Déterminer le domaine de définition d'une fonction racine carrée ou inverse}
\AddCompetence{Fonc21}{1M}{7}{Représenter graphiquement une fonction racine carrée}
\AddCompetence{Fonc22}{1M}{7}{Représenter graphiquement une fonction inverse}
